\section{Background}
\label{sec:background}

This section situates the paper at the intersection of multi-human AI support and team information-processing support. It motivates why a dashboard-centered command-support system can be understood as one instance of a broader class of digital support systems for shared coordination, and why task-state monitoring is a meaningful bounded capability within that class.

A large share of current generative AI applications remains focused on the individual user: question answering, drafting, summarisation, or recommendation for a single human operator \citep{Maedche2019,Nah2023}. By contrast, many work settings require models to interpret contributions distributed across several people, maintain state over time, and support shared rather than individual understanding. Reviews of human-AI teaming show that much of the literature still focuses on bilateral assistants, dyads, or otherwise narrowly scoped team configurations \citep{Lyons2021,ONeill2022,Anthony2023,Bankins2024}. This leaves open how AI systems should support shared state and coordination across ongoing multi-party interaction through a shared artefact such as a dashboard.

Research on Group Support Systems has long shown that collaboration quality depends not only on who contributes, but also on how information is processed and coordinated within the group \citep{Nunamaker1991,DennisValacich1993,DennisWixomVandenberg2001,Briggs2003}. Facilitation has been central in this tradition because groups systematically suffer from process losses such as production blocking, dominance effects, incomplete information pooling, and premature convergence \citep{StasserTitus1985,DennisValacich1993,Briggs2003}. Recent work on AI in teams suggests that digital artefacts may take over selected support capabilities, especially those involving consistency, timing, monitoring, or large-scale information processing, while humans retain judgement, legitimacy, and contextual authority \citep{Dellermann2019HI,Seeber2020,DennisLakhiwalSachdeva2023}.

Our paper focuses on one such capability: maintaining an explicit representation of task state from ongoing multi-party communication. In facilitation terms, this is primarily an information-processing support function. The command-support system does not replace operational decision making. Rather, it externalises on the dashboard what the team has already established: which tasks have been handed out, which unit appears responsible, which tasks have been completed, and what completion result has been reported.

Technically, this formulation is related to dialogue state tracking and information extraction. Dialogue state tracking maintains structured state representations from dialogue history in task-oriented systems \citep{Balaraman2021DSTSurvey}. Information and event extraction similarly recover structured records, events, or arguments from text, often against predefined schemas \citep{ChambersJurafsky2011,XiangWang2019EventSurvey}. Our setting differs in that the state variables are predefined SOP- and procedure-related operational tasks, the evidence is distributed across multi-party radio communication, and the output is used for incremental dashboard support rather than autonomous dialogue management or open-ended event discovery.

We study this capability in a deliberately bounded operational setting rather than in the most open-ended form of multi-human coordination. If a digital support system cannot reliably maintain assignment and completion state for predefined tasks, then more advanced functions such as summarisation, blind-spot detection, or process coaching are unlikely to be trustworthy. Conversely, if task-state monitoring proves robust under reduced transcript structure, this suggests that at least some multi-human information-processing functions may not require expensive preprocessing pipelines.

Taken together, these strands motivate AI support for shared coordination, but they do not yet answer the more specific deployment question studied in this paper: how much transcript structure is required for closed-world task-state monitoring from ongoing multi-party communication? Firefighter incident command is a useful case because the dashboard-centered support function is operationally meaningful, the communication is genuinely multi-party, and transcript enrichment can be costly. For a bounded function such as digital task-state tracking, richer transcript preprocessing may help, but it is not obvious that speaker labels and clean utterance boundaries are necessary.
